% Partie : set-theory
% Chapitre : spine
% Section : idea

\documentclass[../../../include/open-logic-section]{subfiles}

\begin{document}

\olfileid[fr]{sth}{spine}{valpha}

\olsection{Définir les étapes comme les \texorpdfstring{$V_\alpha$}{Vα}}

Dans le \olref[sth][ordinals][]{chap}, nous avons défini les bons ordres
et les ordinaux (de von Neumann). Nous allons maintenant nous en
servir pour caractériser la hiérarchie des ensembles \emph{elle-même}.
Rappelons qu'à la \olref[sth][ordinals][opps]{sec}, nous avons défini
les notions d'ordinal successeur et d'ordinal limite. Elles
interviennent dans la définition suivante :

\begin{defn}\ollabel{defValphas}
	\begin{align*}
	V_\emptyset &\defis \emptyset\\
	V_{\ordsucc{\alpha}} &\defis \Pow{V_\alpha} & &
	\text{pour tout ordinal }\alpha\\
	V_{\alpha} &\defis \bigcup_{\gamma < \alpha} V_\gamma & &
	\text{si }\alpha\text{ est un ordinal limite}
\end{align*}
\end{defn}
\noindent
Il s'agira d'une définition par \emph{récursion transfinie} sur les
ordinaux. Comparons-la aux définitions récursives des fonctions sur
les nombres naturels.\footnote{Voir les définitions de l'addition,
de la multiplication et de l'exponentiation dans la
\olref[sfr][infinite][induction]{sec}. Note de l'édition française : le renvoi original vise la section sur les algèbres de Dedekind ; les équations citées figurent dans la section suivante.} Avec les nombres naturels,
on définit un cas de base et des cas successeurs ; avec les ordinaux,
il faut aussi décrire ce qui se passe aux étapes \emph{limites}.

Cette définition des $V_\alpha$ constituera un jalon important.
Nous avons motivé de manière informelle la hiérarchie des ensembles
en présentant leur formation par \emph{étapes}. Les $V_\alpha$ sont,
en effet, précisément ces étapes. Mais il s'agit, et c'est important,
d'une caractérisation \emph{interne} des étapes. Au lieu de proposer
un éventuel \emph{modèle} de la théorie, nous aurons défini les
étapes \emph{au sein} de notre théorie des ensembles.

\end{document}
